%% ============================================================
%% EXTENDED PHYSICS + VALIDATION PROTOCOL + IMPLEMENTATION -- v2 merge.
%% Drop-in for:  <<KEEP:pressure-drop>> <<KEEP:temp-prop>> <<KEEP:delay>>  (S2.5)
%%               <<KEEP:stage1>> <<KEEP:stage2>>                          (S2.7)
%%               <<KEEP:implementation>>                                  (S2.8)
%% Source: v1 lines 967-1225.
%%
%% THE STRUCTURAL CHANGE. v1 presented L3+ and L3NL as two variants sharing an
%% identical physics scope and differing only in mathematical treatment -- so pressure,
%% temperature propagation and transport delay all switched on together, in one step.
%% That is precisely the confound R1.3 and R2.3 objected to. Here each phenomenon is a
%% separate rung: \Ltwo adds temperature propagation, \Lthree adds trunk pressure and
%% pumping, \Lsix adds transport delay, and \NLref is the nonlinear reference. The
%% mathematics is described per phenomenon, not per variant.
%%
%% NOTE ON \Ltwo: the solved free-temperature MILP is degenerate (the optimiser floors
%% the supply temperature for ~91 % of hours), so temperature propagation is reported
%% through the forward evaluator, not through its own solve. Stated here so the reader
%% meets the caveat in the methods rather than discovering it in the results.
%%
%% NOTE ON \Lsix: at hourly resolution the discretised travel time is zero timesteps for
%% every trunk pipe, so \Lsix reduces to \Lthree identically -- same model size, same
%% solution to the last digit. This is a property of the temporal resolution, and saying
%% so here makes the +0.00 % result in S4.7 self-evident rather than suspicious.
%%
%% Packages: no \si/\SI/\ce -- plain units, matching paper_COMPILE.tex.
%% ============================================================

%% ---------- S2.5 Extended thermo-hydraulic formulation ----------

Three phenomena are added above the node-resolved baseline $\Lone$, each on its own rung
so that no comparison changes more than one thing: temperature propagation along pipes
($\Ltwo$), trunk pressure drop and the pumping power it implies ($\Lthree$), and
transport delay ($\Lsix$). A fourth object, the nonlinear reference $\NLref$, retains all
three in native form and is used for evaluation and for bounded re-solves rather than as a
rung of the ladder. Piecewise-linear approximations use SOS2 machinery with $K = 3$
segments and $K + 1 = 4$ support points throughout.

\paragraph{Pressure drop and pumping power ($\Lthree$).}
Friction in pipe $p$ follows Darcy--Weisbach, relating the pressure drop $\Delta p_p$ [Pa]
to the square of the mass flow $\dot{m}_p$ [kg/s]:
%
\begin{equation}
  \Delta p_{p} = f_D\, \frac{L_p}{D_p}\,
    \frac{8\,\dot{m}_p^2}{\rho \pi^2 D_p^4},
  \label{eq:darcy}
\end{equation}
%
with $f_D$ the Darcy friction factor, $D_p$ [m] the inner diameter and $\rho$ [kg/m$^3$]
the fluid density. Because the supply and return temperatures are prescribed by the
heating curve, the mass flow is linear in the pipe heat flow $Q_{p,t}$ [MW]:
%
\begin{equation}
  \dot{m}_{p,t} = \frac{Q_{p,t}}{
    c_p \bigl(T_{\mathrm{sup}}^{\mathrm{nom}}(t)
    - T_{\mathrm{ret}}^{\mathrm{nom}}(t)\bigr)},
  \label{eq:massflow}
\end{equation}
%
so that $\Delta p_{p,t} = R_p\, Q_{p,t}^2$ with $R_p$ [Pa/MW$^2$] the hydraulic
resistance. In $\Lthree$ this quadratic is approximated piecewise-linearly over $K$ equal
segments with SOS2 weights $w_{p,t,k} \in [0,1]$,
%
\begin{equation}
  \Delta p_{p,t}^{\mathrm{PWL}} = \sum_{k=1}^{K}
    m_k^{(p)}\, w_{p,t,k},
  \label{eq:pressure_pwl}
\end{equation}
%
whose maximum approximation error is
$\varepsilon^{\max}_{\Delta p} = R_p \bar{Q}_p^2 / (4K^2)$, about $2.8$\,\% of the peak
pressure drop at $K = 3$ (Appendix~\ref{app:pwl}). In $\NLref$ the quadratic enters
natively.

Pumping power is the hydraulic work summed over pipes and divided by the pump efficiency,
%
\begin{equation}
  P_t^{\mathrm{pump}} = \frac{1}{\eta_{\mathrm{pump}}}
    \sum_{p \in \mathcal{P}}
    \frac{\dot{m}_{p,t}\, \Delta p_{p,t}}{\rho},
  \label{eq:pump_power}
\end{equation}
%
which is cubic in $Q_{p,t}$ and is piecewise-linearised directly in $\Lthree$; in $\NLref$
it is carried by the auxiliary $W_{p,t} = Q_{p,t}^2$ entering as $Q_{p,t} W_{p,t}$. The
pump efficiency is taken from the installed units' datasheets rather than assumed. Both
the supply and the return circuit are resolved, and every consumer additionally requires a
differential pressure across its substation, set to the operator's specified $0.6$\,bar --
the term omitted in the original submission.

\paragraph{Temperature propagation ($\Ltwo$).}
Where $\Lone$ computes pipe losses from fixed nominal temperatures, $\Ltwo$ tracks the
outlet temperature as the fluid travels. In steady state it decays exponentially from the
inlet temperature towards the ground temperature $T_{\mathrm{gr}}$:
%
\begin{equation}
  T_{p,t}^{\mathrm{out}} = T_{\mathrm{gr}}(t)
    + \bigl(T_{n_1,t}^{\mathrm{sup}} - T_{\mathrm{gr}}(t)\bigr)
    \cdot \underbrace{\exp\!\left(
      \frac{-U_p L_p}{\dot{m}_{p,t}\, c_p}
    \right)}_{\displaystyle\phi_p(\dot{m}_{p,t})}.
  \label{eq:temp_prop}
\end{equation}
%
The decay factor $\phi_p \in (0,1]$ is the fraction of the inlet temperature excess
surviving to the outlet: near unity at high flow, near zero as flow approaches zero. Two
nonlinearities follow -- the exponential itself and the bilinear product
$T_{n_1,t}^{\mathrm{sup}} \phi_p$. The exponential is approximated by a PWL in $\dot{m}$
shared by every level that uses it, so that this component of the error is common and
cancels in the comparisons. The bilinear product enters natively in $\NLref$; in $\Ltwo$
it is linearised by a first-order Taylor expansion about the nominal operating point,
%
\begin{align}
  T_{p,t}^{\mathrm{out}} &\approx T_{\mathrm{gr}}(t)
    + \bigl(T^{\mathrm{nom}}(t) - T_{\mathrm{gr}}(t)\bigr)
      \phi_{p,t}^{\mathrm{PWL}} \notag \\
    &\quad + \phi_p^{\mathrm{nom}}
      \bigl(T_{n_1,t}^{\mathrm{sup}} - T^{\mathrm{nom}}(t)\bigr),
  \label{eq:temp_prop_lin}
\end{align}
%
with the neglected second-order cross-term bounded by
$|\Delta T_{\max}||\Delta\phi_{\max}| < 2.5$\,K (Appendix~\ref{app:taylor_derivation}).

One property of this formulation determines how we report it. Freeing the node
temperatures makes the programme non-convex and, absent a hydraulic penalty on the
resulting flow increase, degenerate: the optimiser reduces the supply temperature
uniformly to its lower bound, collapsing the loss term rather than reproducing the
physical decay. In our solve it does so for about $91$\,\% of hours, against a mean
supply temperature of $85.9$\,°C at the neighbouring levels. The solved objective of
$\Ltwo$ is therefore not a meaningful fidelity measurement, and we do not use it. The
effect of temperature propagation is instead quantified through the forward evaluator of
Section~\ref{subsec:evaluator}, which applies the native exponential to a fixed schedule
and cannot exploit the degeneracy. $\Ltwo$ remains in the ladder as the formulation whose
linearisation error is being measured, not as a source of dispatch cost.

\paragraph{Transport delay ($\Lsix$).}
At finite velocity, heat injected at the source at time $t$ arrives downstream later. The
nominal travel time through pipe $p$, with cross-section $A_p = \pi D_p^2/4$, is
%
\begin{equation}
  \tau_p = \frac{L_p\, \rho\, A_p}{\dot{m}_p^{\mathrm{nom}}},
  \label{eq:delay_time}
\end{equation}
%
discretised to $k_p = \mathrm{round}(\tau_p/\Delta t)$ time steps, so that the heat
delivered downstream is the heat injected $k_p$ steps earlier less the corresponding
losses:
%
\begin{equation}
  Q_{p,t}^{\mathrm{delivered}} =
    Q_{p,\,t-k_p} - Q_{p,\,t-k_p}^{\mathrm{loss,pipe}}
  \quad \forall\, p,\; t \geq k_p{+}1.
  \label{eq:delayed_balance}
\end{equation}
%
Isolating delay on its own rung is what allows it to be separated from the linearisation
error, which the original submission could not do. It also produces a clean negative
result, and the reason is visible in \eqref{eq:delay_time}: on this network every trunk
travel time is shorter than an hour, so $k_p = 0$ for every pipe and
\eqref{eq:delayed_balance} collapses to the undelayed balance. $\Lsix$ and $\Lthree$ are
consequently the same programme -- identical variable, binary and constraint counts, and
an identical optimum to the last digit. The delay effect at hourly resolution is exactly
zero, not merely small, and this is a statement about the temporal resolution rather than
about transport delay in general: at sub-hourly resolution, or on a network with
kilometres more trunk, $k_p$ would not vanish.

%% ---------- S2.7 Validation protocol ----------

The network was upgraded with the heat pump, electrode boiler and thermal storage
\emph{after} the measurement record was collected, so validation follows a two-stage
strategy after \citet{Kus2025} and \citet{Maldonado2024}.

\paragraph{Stage 1: direct network validation.}
The pipe model is validated against pre-upgrade monitoring data with the new assets set to
zero capacity, calibrated by the branch-sequential $U_p$-tuning procedure of
\citet{Maldonado2024} but constrained, in this revision, to a defensible coefficient range
(Section~\ref{subsec:calib}). Table~\ref{tab:val_targets} lists the key performance
indicators and their gates, all defined \emph{ex ante} from published practice and fixed
before the calibration was run.

We draw attention to how these gates are used, because it differs from the original
submission. They are reported in full in Section~\ref{subsec:validation}, including the
several that are not met. The temperature gates in particular cannot be met by any model
of this network: the consumer sensors are billing instruments sited downstream of
three-way mixing valves, so the metered temperature sits systematically below the primary
junction temperature, and a point-in-time match is additionally sensitive to a network flow
the billing metering does not pin down. Retaining gates that the instrumentation cannot
support, and reporting the failures, is the point rather than an oversight: it is the
evidence for the claim that a model resolved finer than its metering cannot have that
resolution verified. The conclusions are therefore anchored to the annual delivered energy
-- and hence the annual network loss -- which is both gated and met, and which is the
quantity the cost results actually depend on.

\paragraph{Stage 2: necessary-condition verification.}
Dispatch of the new assets is verified against plausibility bounds rather than
measurements: heat-pump COP within $[2.5, 5.5]$, thermal-storage state of charge within
$[0.05, 0.95]\bar{E}_s$, no simultaneous charge and discharge, and per-step energy
closure. These are necessary conditions, not sufficient ones, and we say so; a
post-upgrade measurement campaign is the natural follow-up and is listed among the
limitations.

\paragraph{Structural sanity checks.}
Two checks confirm internal consistency. First, the nonlinear reference with its quadratic
constraints deactivated reproduces the baseline objective to within $0.01$\,\%. Second, the
relaxation property $\mathrm{cost}(\CP) \le \mathrm{cost}(\Lone)$ holds on every one of the
135 synthetic instances -- a copperplate can never cost more than the resolved model of the
same network, so a violation would indicate a formulation error rather than a finding.

%% ---------- S2.8 Implementation ----------

All levels are implemented in Python with Pyomo~\cite{pyomo}, instantiated from a single
parameterised codebase through YAML configuration files, so that a level is a
configuration rather than a separate model; this is what makes the one-phenomenon-per-step
contrasts exact. Solves use Gurobi with identical random seeds on a 66-core, 180\,GB
machine.

Optimality tolerances are tighter than in the original submission, and deliberately so.
The decomposition terms of Section~\ref{subsec:decomp} are small enough that the original
$0.5$\,\% tolerance placed the topology and interaction terms inside solver noise -- at
that tolerance the interaction carried the opposite sign to its converged value. All
linear levels are therefore solved to a $0.01$\,\% gap with \texttt{MIPFocus=2}, and the
synthetic factorial to $0.1$\,\% or better, so that every reported term sits well above
the attained gap. Where a bilinear formulation is involved the tolerance is stated
alongside the effect it measures, and no effect smaller than the attained tolerance is
claimed.

The station-resolved levels $\Lfour$ and $\Lfive$ are evaluated by the validated forward
pressure-study module rather than re-solved as dispatch problems. Forward evaluation
upper-bounds their regret, so the null they establish is rigorous rather than assumed, and
this keeps the solved baseline byte-identical across the ladder. The same applies to
$\Ltwo$, for the degeneracy reason given above.
